\ysrc{251}
由猜想 6.1.15 (2) 立得 $\mathfrak A\bigl((c\text{-}\mathrm{Ind}_K^G V)^\wedge\bigr)$ 集中在 0 次。事实上，由注 6.1.24 有
\[
\mathfrak A(c\text{-}\mathrm{Ind}_K^G V)
=L_\infty\widehat{\otimes}^L_{\mathcal O[[G]]}(c\text{-}\mathrm{Ind}_K^G V)^\wedge
=L_\infty\widehat{\otimes}^L_{\mathcal O[[K]]}V
\]
（其中 $\widehat{\otimes}$ 表示 $p$-进完备张量积），于是只要假定 $L_\infty$ 在 $\mathcal O[[K]]$ 上平坦，结论即得。

特别地，性质 (4) 中考虑的那些层 $\mathfrak A\bigl(\widehat{\Pi(\underline\lambda,\tau)^\circ}\bigr)$ 都集中在 0 次。

\begin{rem}{6.1.28}
关于猜想第 (1) 条所猜的有界振幅，我们其实可以说得精确得多：我们预期函子 $\mathfrak A$ 的振幅为 $[1-d,0]$。

事实上，当 $d=2$、$p>2$ 时，这一点可由注 6.1.27 以及注 7.3.11 中解释的消解得出（至少在下述意义下成立：我们忽略中心的作用，因为不预期它对振幅有贡献）；并且我们预期对一切 $d$ 都存在类似的消解。
\end{rem}

\begin{rem}{6.1.29}
按构造，$\mathrm{Ind}\,D^b_{\mathrm{f.p.}}(\mathrm{sm}\,G)$ 中的紧对象恰是 $D^b_{\mathrm{f.p.}}(\mathrm{sm}\,G)$ 的对象，而 $\mathrm{Ind}\,\mathrm{Coh}(\mathcal X_d)$ 中的紧对象由 $D^b_{\mathrm{coh}}(\mathcal X_d)$ 给出。于是由引理 A.4.6 (3) 可知：若 $\pi$ 是 $\mathrm{Ind}\,D^b_{\mathrm{f.p.}}(\mathrm{sm}\,G)$ 的对象，则 $\pi$ 实际属于 $D^b_{\mathrm{f.p.}}(\mathrm{sm}\,G)$，当且仅当 $\mathfrak A(\pi)$ 是 $D^b_{\mathrm{coh}}(\mathcal X_d)$ 的对象。

特别地，若 $\pi$ 集中在 0 次，则 $\pi$ 有限表现当且仅当 $\mathfrak A(\pi)$ 是上同调层凝聚的有界复形。（事实上由注 6.1.28，我们预期 $\mathfrak A(\pi)$ 自动有界。）

若 $\pi$ 集中在 0 次、有限生成而非有限表现，则 $\mathfrak A(\pi)$ 不凝聚。不过我们断言 $H^0(\mathfrak A(\pi))$ 凝聚。事实上，$\pi$ 既有限生成，就存在某个有限长 $K$-表示 $U$ 以及满射 $c\text{-}\mathrm{Ind}_K^G U\to\pi$。由于 $\mathfrak A$ 右 $t$-正合（见注 6.1.28），得到满射 $H^0\bigl(\mathfrak A(c\text{-}\mathrm{Ind}_K^G U)\bigr)\to H^0(\mathfrak A(\pi))$；而 $\mathfrak A(c\text{-}\mathrm{Ind}_K^G U)$ 凝聚（因为 $c\text{-}\mathrm{Ind}_K^G U$ 有限表现）。
\end{rem}

\begin{rem}{6.1.30}
我们曾力求只用一组极少的、有用的性质来表述猜想 6.1.15；下面诸注解释这些性质的部分动机。这些性质究竟在多大程度上唯一确定 $\mathfrak A$，目前并不清楚，但在实践中它们似乎对它构成了相当强的约束。
\end{rem}

\begin{rem}{6.1.31}
可以把 $L_\infty$ 看作完备上同调的一个「泛的」、纯局部的版本。更确切地说，在假定 $p\nmid 2d$ 之下，对任意一点 $x\in\mathcal X_d(\overline{\F}_p)$，文 [CEG$^+$16] 在 $x$ 处构造了一个通用（versal）态射 $f:\Spf R_\infty\to\mathcal X_d$，并配以一个带有 $G$ 的交换作用的 $R_\infty$-模 $M_\infty$。这里 $R_\infty$ 是某个幂级数环，它建立在对应于 $x$ 的 $\Gal_F$-表示的泛提升环之上，其中的幂级数变量正是 Taylor--Wiles--Kisin 方法中出现的那组「拼接变量」；而 $M_\infty$ 则由对某些酉群的完备上同调施行拼接方法得到。（模曲线的类似构造见 §3，两者极为相似。）

$R_\infty$ 与 $M_\infty$ 的构造是整体性的，且依赖于许多选择；但我们预期 $M_\infty=f^*L_\infty$（另见预期定理 9.4.2）。这将证实 [CEG$^+$16, §6] 中的猜想，即 $M_\infty$ 应当是纯局部的；事实上，该文所假设的 $L_\infty$ 不过是我们这个 $L_\infty$ 到局部形变环上的拉回。\ysrc{252}这也解释了我们关于 $L_\infty$ 的一部分预期；例如 $M_\infty$ 按构造在 $\mathcal O[[K]]$ 上平坦，由此便得到「$L_\infty$ 在 $\mathcal O[[K]]$ 上平坦」这一预期。
\end{rem}

\begin{rem}{6.1.32}
猜想 6.1.15 的性质 (2) 是下述预期的一个类比：[BZCHN24, Conj. 4.15] 断言凝聚 Springer 层集中在 0 次。在我们的情形，它由注 6.1.31 的预期所驱动（该预期蕴含：$L_\infty$ 到任意通用环的拉回都集中在 0 次）。
\end{rem}

\begin{rem}{6.1.33}
由性质 (4)，对任意正则 Hodge 型 $\underline\lambda$ 与惯性型 $\tau$，层 $\mathfrak A\bigl(\widehat{\Pi(\underline\lambda,\tau)^\circ}\bigr)$ 支撑在 $\mathcal X_d^{\mathrm{crys},\underline\lambda,\tau}$ 上。我们有闭浸入 $i:\mathcal X_d^{\mathrm{crys},\underline\lambda,\tau}\hookrightarrow\mathcal X_d$，并预期它是纯的、余维为 $[F:\Q_p]d(d+1)/2$（我们之所以说「预期」，是因为文献中尚无形式代数叠的维数理论；但注意由 [Kis10b, BIP23a] 已知相应的论断在通用环上成立，即在伽罗瓦形变环上成立）。承认这一预期，便有 $i^!\omega_{\mathcal X_d}=\omega_{\mathcal X_d^{\mathrm{crys},\underline\lambda,\tau}}\bigl[-[F:\Q_p]d(d+1)/2\bigr]$，其中我们滥用记号，用 $\omega_{\mathcal X_d^{\mathrm{crys},\underline\lambda,\tau}}$ 表示按 $\mathcal X_d^{\mathrm{crys},\underline\lambda,\tau}$ 的维数作平移后的对偶复形（于是若 $\mathcal X_d^{\mathrm{crys},\underline\lambda,\tau}$ 是 Cohen--Macaulay 的，则 $\omega_{\mathcal X_d^{\mathrm{crys},\underline\lambda,\tau}}$ 是集中在 0 次的一个层）。因此对任意 $\mathcal F\in D_{\mathrm{Coh}}(\mathcal X_d^{\mathrm{crys},\underline\lambda,\tau})$，有
\[
\begin{aligned}
&i_*\mathrm{R}\mathcal Hom_{\mathcal O_{\mathcal X_d^{\mathrm{crys},\underline\lambda,\tau}}}
\bigl(\mathcal F,\omega_{\mathcal X_d^{\mathrm{crys},\underline\lambda,\tau}}\bigr)\\
&\qquad=\mathrm{R}\mathcal Hom_{\mathcal O_{\mathcal X_d}}\bigl(i_*\mathcal F,\omega_{\mathcal X_d}\bigr)
\bigl[[F:\Q_p]d(d+1)/2\bigr].
\end{aligned}
\]
记
\[
\mathrm{D}_{\mathcal O_{\mathcal X_d^{\mathrm{crys},\underline\lambda,\tau}}}:
D_{\mathrm{Coh}}\bigl(\mathcal X_d^{\mathrm{crys},\underline\lambda,\tau}\bigr)
\to D_{\mathrm{Coh}}\bigl(\mathcal X_d^{\mathrm{crys},-\underline\lambda,\tau^\vee}\bigr)
\]
为 Grothendieck--Serre 对偶与 $i^*$ 的复合，其中
\[
\iota:\mathcal X_d^{\mathrm{crys},-\underline\lambda,\tau^\vee}
\xrightarrow{\ \sim\ }\mathcal X_d^{\mathrm{crys},\underline\lambda,\tau}
\]
由 $\mathcal X_d$ 的对合 $\iota$ 诱导。

取 $\mathcal F=\mathfrak A\bigl(\widehat{\Pi(\underline\lambda,\tau)^\circ}\bigr)$。注意由 (6.1.7) 可得
\[
\begin{aligned}
\mathbf D\bigl(\widehat{\Pi(\underline\lambda,\tau)^\circ}\bigr)
&:=\mathbf D\Bigl(\bigl(c\text{-}\mathrm{Ind}_K^G(W_{\underline\lambda}
\otimes_{\mathcal O}\sigma^{\mathrm{crys},\circ}(\tau))\bigr)^\wedge\Bigr)\\
&=\Bigl(\bigl(c\text{-}\mathrm{Ind}_K^G(W^*_{-\underline\lambda}
\otimes_{\mathcal O}\sigma^{\mathrm{crys},\circ}(\tau^\vee))\bigr)^\wedge\Bigr)[1],
\end{aligned}
\]
其中我们滥用记号，用 $W^*_{-\underline\lambda}$ 表示那个（并非对偶的！）Weyl 模，并取 $\sigma^{\mathrm{crys},\circ}(\tau^\vee)$ 为 $\sigma^{\mathrm{crys}}(\tau^\vee)$ 中与我们选定的格 $\sigma^{\mathrm{crys},\circ}(\tau)$ 相对偶的格。顾及 (6.1.16)，便得到
\begin{equation}\tag{6.1.34}\label{eq:a4-6-1-34}
\begin{aligned}
&\mathfrak A\Bigl(\bigl(c\text{-}\mathrm{Ind}_K^G(W^*_{-\underline\lambda}
\otimes_{\mathcal O}\sigma^{\mathrm{crys},\circ}(\tau^\vee))\bigr)^\wedge\Bigr)\\
&\qquad\xrightarrow{\ \sim\ }
\mathrm{D}_{\mathcal O_{\mathcal X_d^{\mathrm{crys},\underline\lambda,\tau}}}
\bigl(\widehat{\Pi(\underline\lambda,\tau)^\circ}\bigr).
\end{aligned}
\end{equation}
\end{rem}

\begin{rem}{6.1.35}
猜想 6.1.15 的性质 (3) 可以看作在解释下述事实：那些拼接模（由有限层次上的 Taylor--Wiles--Kisin 拼接构造）在各自的支撑上是极大 Cohen--Macaulay 的。特别地，我们预期可以把性质 (4) 加强为：对任意正则 Hodge 型 $\underline\lambda$ 与惯性型 $\tau$，$\mathfrak A\bigl(\widehat{\Pi(\underline\lambda,\tau)^\circ}\bigr)$ 集中在 0 次，在 $\mathcal X_d^{\mathrm{crys},\underline\lambda,\tau}$ 上极大 Cohen--Macaulay，并且 $\mathfrak A\bigl(\widehat{\Pi(\underline\lambda,\tau)^\circ}\bigr)^{\mathrm{rig}}_\eta$ 在 $(\mathcal X_d^{\mathrm{crys},\underline\lambda,\tau})^{\mathrm{rig}}_\eta$ 上是秩一的局部自由。下面我们勾勒这一点的解释。
\end{rem}

\ysrc{253}
既然我们假定 $\mathfrak A\bigl(\widehat{\Pi(\underline\lambda,\tau)^\circ}\bigr)$ 集中在 0 次且支撑在整个 $\mathcal X_d^{\mathrm{crys},\underline\lambda,\tau}$ 上，而 $\mathfrak A\bigl(\bigl(c\text{-}\mathrm{Ind}_K^G(W^*_{-\underline\lambda}\otimes_{\mathcal O}\sigma^{\mathrm{crys},\circ}(\tau^\vee))\bigr)^\wedge\bigr)$ 亦然\footnote{我们并未就对偶 Weyl 模 $W^*_{-\underline\lambda}$ 写出显式的猜想，但这一论断可由关于 $W_{-\underline\lambda}$ 的版本推出，因为这些表示在把 $p$ 可逆化之后是同构的。}，则由 (6.1.34) 可知 $\mathfrak A\bigl(\widehat{\Pi(\underline\lambda,\tau)^\circ}\bigr)$ 在 $\mathcal X_d^{\mathrm{crys},\underline\lambda,\tau}$ 上极大 Cohen--Macaulay（见 [Sta, Tag 0B5A]）。

由此可知 $\mathfrak A\bigl(\widehat{\Pi(\underline\lambda,\tau)^\circ}\bigr)^{\mathrm{rig}}_\eta$ 在 $(\mathcal X_d^{\mathrm{crys},\underline\lambda,\tau})^{\mathrm{rig}}_\eta$ 上极大 Cohen--Macaulay。由于 $(\mathcal X_d^{\mathrm{crys},\underline\lambda,\tau})^{\mathrm{rig}}_\eta$ 光滑，这说明 $\mathfrak A\bigl(\widehat{\Pi(\underline\lambda,\tau)^\circ}\bigr)^{\mathrm{rig}}_\eta$ 局部自由；并且我们预期 $\mathfrak A$ 的全忠实性连同 [CEG$^+$16, Cor. 3.12] 还蕴含它是秩一的局部自由。

回忆一下，拼接模的 Cohen--Macaulay 性质通常是用完全不同的手段建立的，即利用钻石算子所构成的环 $S_\infty$ 的作用；见 [Dia97]。

\begin{rem}{6.1.36}
正如在定义 $\mathcal X_d$ 上 Grothendieck--Serre 对偶的修正版本 $\mathbf D_{\mathcal X_d}$ 时已经指出的，$\mathcal X_d$ 上有「主对合」$\iota$，定义为 $\rho\mapsto\rho^\vee$。它在 $D^b_{\mathrm{coh}}(\mathcal X_d)$ 上诱导一个自对合 $\iota^*$。另一方面，$\mathrm{GL}_d(F)$ 上有通常的主对合 $g\mapsto(g^{-1})^t$，用它对 $\mathrm{GL}_d(F)$ 的表示作预复合，便在 $D^b_{\mathrm{f.p.}}(\mathrm{sm}\,G)$ 上诱导一个自对合 $\iota'$。我们预期 $\mathfrak A$ 把对合 $\iota'$ 与 $\iota$ 相缠。

例如，在猜想 6.1.15 的性质 (3) 的陈述中，与其把 $\mathcal X_d$ 上的 Grothendieck--Serre 对偶 $\mathbf D$ 换成它的扭 $\mathbf D_{\mathcal X_d}:=\iota^*\mathbf D$，我们也可以把 $D^b_{\mathrm{f.p.}}(\mathrm{sm}\,G)$ 上的对偶 $\mathbf D$ 换成它的扭 $\mathbf D':=\iota'\circ\mathbf D$。于是性质 (3) 可以改述为：$\mathfrak A$ 互换 $\mathbf D'$ 与 $\mathbf D$（差一个平移）。此时 (6.1.7) 蕴含（忽略可能的平移）：若 $V$ 是 $K$ 的一个不可约模 $p$ 表示，则 $\mathbf D'\bigl((c\text{-}\mathrm{Ind}_K^G V)^\wedge\bigr)=\bigl(c\text{-}\mathrm{Ind}_K^G V^\circ\bigr)^\wedge$；若 $V^\circ$ 是不可约局部代数表示 $V$ 的一个 $K$-不变格，则 $\mathbf D'\bigl((c\text{-}\mathrm{Ind}_K^G V^\circ)^\wedge\bigr)=\bigl(c\text{-}\mathrm{Ind}_K^G V^\circ\bigr)^\wedge$。因此我们预期相应的凝聚层 $\mathfrak A(c\text{-}\mathrm{Ind}_K^G V)$（相应地 $\mathfrak A\bigl((c\text{-}\mathrm{Ind}_K^G V^\circ)^\wedge\bigr)$）是 Grothendieck--Serre 自对偶的（差一个平移）。
\end{rem}

\begin{rem}{6.1.37}
我们还预期与抛物诱导相容，正如 $\ell\neq p$ 情形与解析情形那样（后者见 §6.2.14）。
\end{rem}

\subsubsection*{6.1.38.\quad 几何 Breuil--Mézard 猜想}

叠 $\mathcal X_d$ 的几何 Breuil--Mézard 猜想是在 [EG23] 中表述的；它的动机来自（并且实际上等价于）[BM14, EG14] 中关于伽罗瓦形变环的较早的几何 Breuil--Mézard 猜想。不熟悉 Breuil--Mézard 猜想的读者可参阅这些文献（及其中的参考文献），以了解该猜想最初的动机，它们来自自守提升定理。本节说明猜想 6.1.15 如何给出 Breuil--Mézard 猜想的一个进一步加细（把该陈述从闭链的相等提升为凝聚层的 Grothendieck 群中的相等）。

所谓 Serre 权 $\underline k$，是指一组整数 $\{k_{\bar\sigma,i}\}_{\bar\sigma:k\hookrightarrow\overline{\F}_p,\,1\le i\le d}$，满足
\begin{itemize}
\item 对每个 $1\le i\le d-1$ 有 $p-1\ge k_{\bar\sigma,i}-k_{\bar\sigma,i+1}\ge 0$；并且
\item $p-1\ge k_{\bar\sigma,d}\ge 0$，且并非所有 $k_{\bar\sigma,d}$ 都等于 $p-1$。
\end{itemize}

\ysrc{254}
Serre 权的集合与 $\mathrm{GL}_d(k_F)$ 的不可约 $\overline{\F}_p$-表示的集合一一对应，对应方式是取最高权向量（例如见 [Her09] 的附录）；对每个 Serre 权 $\underline k$，我们用 $F_{\underline k}$ 表示 $\mathrm{GL}_d(k_F)$ 相应的不可约 $k$-表示。

用 $\overline{\sigma}^{\mathrm{crys}}(\underline\lambda,\tau)$ 表示由 $W_{\underline\lambda}\otimes_{\mathcal O}\sigma^{\mathrm{crys},\circ}(\tau)\otimes_{\mathcal O}k$ 给出的 $\mathrm{GL}_d(k_F)$ 的 $k$-表示的半单化。存在唯一的整数 $n^{\mathrm{crys}}_{\underline k}(\underline\lambda,\tau)$，使得
\[
\overline{\sigma}^{\mathrm{crys}}(\underline\lambda,\tau)
\cong\oplus_{\underline k}F_{\underline k}^{\oplus n^{\mathrm{crys}}_{\underline k}(\underline\lambda,\tau)}.
\]
用 $Z_{\mathrm{crys},\underline\lambda,\tau}$ 表示对应于 $\mathcal X_d^{\mathrm{crys},\underline\lambda,\tau}$ 的特殊纤维的闭链（即自由阿贝尔群 $\mathbf Z[\mathcal X_{d,\mathrm{red}}]$ 中的一个元素，该群以 $\mathcal X_{d,\mathrm{red}}$ 的不可约分支为基）。下述命题即 [EG23, Conj. 8.2.2]。

\begin{conj}{6.1.39}\label{conj:a4-6-1-39}
存在闭链 $Z_{\underline k}$，使得对每个正则 Hodge 型 $\underline\lambda$ 与每个惯性型 $\tau$，都有 $Z_{\mathrm{crys},\underline\lambda,\tau}=\sum_{\underline k}n^{\mathrm{crys}}_{\underline k}(\underline\lambda,\tau)\cdot Z_{\underline k}$。
\end{conj}

假定猜想 6.1.15 成立。对每个正则 Hodge 型 $\underline\lambda$ 与每个惯性型 $\tau$，在如上选定 $\sigma^{\mathrm{crys}}(\tau)$ 中的一个格 $\sigma^{\mathrm{crys},\circ}(\tau)$ 之后\footnote{由于我们过渡到相应的闭链，格的具体选法无关紧要。}，并（同样如上）置 $\Pi(\underline\lambda,\tau)^\circ:=c\text{-}\mathrm{Ind}_K^G\bigl(W_{\underline\lambda}\otimes_{\mathcal O}\sigma^{\mathrm{crys},\circ}(\tau)\bigr)$，我们令
\[
\mathcal F_{\underline\lambda,\tau}=\mathfrak A\bigl(\widehat{\Pi(\underline\lambda,\tau)^\circ}\bigr),
\]
由注 6.1.27，它是 $\mathcal X_d$ 上的一个凝聚层（集中在 0 次）。对每个 Serre 权 $\underline k$，令
\[
\mathcal F_{\underline k}:=\mathfrak A\bigl(c\text{-}\mathrm{Ind}_K^G F_{\underline k}\bigr).
\]
用 [EG14, Lem. 4.1.1] 作一次容易的归纳即可看出，$\mathcal F_{\underline k}$ 是 $\mathcal X_d$ 的特殊纤维上的一个凝聚层（同样集中在 0 次）。

既然 $\mathfrak A$ 是函子，由定义即知在 $K_0\bigl(\mathrm{Coh}(\mathcal X_d)\bigr)$ 中有等式
\begin{equation}\tag{6.1.40}\label{eq:a4-6-1-40}
[\mathcal F_{\underline\lambda,\tau}\otimes_{\mathcal O}k]
=\sum_{\underline k}n^{\mathrm{crys}}_{\underline k}(\underline\lambda,\tau)\cdot[\mathcal F_{\underline k}].
\end{equation}

正如注 6.1.35 所解释的，$\mathcal F_{\underline\lambda,\tau}$ 在其支撑 $\mathcal X_d^{\mathrm{crys},\underline\lambda,\tau}$ 上极大 Cohen--Macaulay，因此 $\mathcal F_{\underline\lambda,\tau}\otimes_{\mathcal O}k$ 的支撑等于 $Z_{\mathrm{crys},\underline\lambda,\tau}$。再用 [EG14, Lem. 4.1.1] 作一次容易的归纳可知，每个 $\mathcal F_{\underline k}$ 的支撑都是 $\mathbf Z[\mathcal X_{d,\mathrm{red}}]$ 中的一个闭链 $Z_{\underline k}$，于是猜想 6.1.39 是 (6.1.40) 的直接推论。

\subsection{预期：解析情形}

本节中令 $G=\mathrm{GL}_d(F)$，并把它看作 $\Q_p$-解析群。此外，本节假定所固定的 $\Q_p$ 的有限扩张 $L$ 包含 $F$ 的正规闭包。我们想表述一个与猜想 6.1.15 平行的猜想，其中涉及局部解析表示以及 $\mathcal X_d$ 上的层。为此，须先讨论局部解析表示语境下的导出范畴。

\ysrc{255}
\subsubsection*{6.2.1.\quad 局部解析表示的导出范畴}

用 $\mathcal D(G)$ 表示 $G$ 的系数在 $L$ 中的分布代数。该分布代数定义为局部解析函数空间 $G\to L$ 的对偶空间，精确定义见 [ST02]（在我们的情形，系数域总是默认已知的，故从记号中略去）。

由 Schneider--Teitelbaum 的一个定理 [ST02, Corollary 3.3]，$G$ 在紧型 $L$-向量空间上的局部解析表示连同连续 $G$-态射所构成的范畴，反等价于核 Fréchet 空间上的分别连续 $\mathcal D(G)$-模连同连续 $\mathcal D(G)$-模同态所构成的范畴。这一反等价把紧型 $L$-向量空间 $V$ 送到它的强对偶 $V'=V'_b$。在这一范畴反等价之下，容许局部解析 $G$-表示（按定义）对应于余容许 $\mathcal D(G)$-模，即满足下述条件的那些 $\mathcal D(G)$-模 $M$：$M$ 在 $G$ 的极大紧子群 $K=\mathrm{GL}_n(\mathcal O_F)$\tn{底本此处印作 $n$，而上文取 $G=\mathrm{GL}_d$，$n$ 并未定义；据上下文当为 $d$。译文照底本印作 $n$。} 的分布代数 $\mathcal D(K)$ 上余容许。更确切地说，这意味着 $M$ 是 $\mathcal D(K)$ 的某些 Banach 完备化 $\mathcal D_r(K)$ 上有限型模的逆向极限，参见 [ST03, §3, 4]。下文中我们用 $\mathrm{Rep}^{\mathrm{an}}_L G$ 表示紧型 $L$-向量空间上局部解析 $G$-表示所构成的范畴。

我们还要指出，局部解析表示 $\pi$ 上的 $G$-作用自动延拓为分布代数 $\mathcal D(G)$ 的分别连续作用，即 $\pi$ 可以看作一个 $\mathcal D(G)$-模 [ST02, Prop. 3.2]。与范畴 $D(\mathrm{sm}\,G)$ 类似，我们考虑局部解析表示的复形所构成的导出范畴 $D(\mathrm{an}\,G)$。该范畴的定义（以及从它到其他导出范畴的导出函子的定义）会遇到一些困难，因为我们是在同调代数的语境中与拓扑对象打交道。这类问题最好借助 Clausen--Scholze 引入的凝聚态（condensed）结构与固态模来处理；而固态模框架下的局部解析表示理论由 Rodrigues Jacinto 与 Rodríguez Camargo 在 [RJRC22] 与 [RJRC25] 中发展起来。我们简述\footnote{不熟悉固态模形式体系的读者可以把范畴 $D(\mathrm{an}\,G)$ 当作一个黑箱，只把它看作 $\mathrm{Rep}^{\mathrm{an}}_L G$ 在适当意义下的导出范畴。} $(\infty)$-范畴 $D(\mathrm{an}\,G)$ 的构造：在 $L^\blacksquare$-模的导出范畴中，$L^\blacksquare[G]$-模所构成的 $(\infty)$-范畴 $D_{L^\blacksquare}(L^\blacksquare[G])$ 有一个全子范畴 $D_{L^\blacksquare}(\mathcal D(G))$，由该范畴中的 $\mathcal D(G)$-模组成，见 [RJRC22, Cor. 4.3.4] 与 [RJRC25, Remark 1.2.1 (2)]。范畴 $D_{L^\blacksquare}(\mathcal D(G))$ 带有一个自函子
\[
(-)^{\mathrm{Rla}}:D_{L^\blacksquare}(\mathcal D(G))\longrightarrow D_{L^\blacksquare}(\mathcal D(G)),
\]
它把（复形的）$\mathcal D(G)$-模送到其（导出）局部解析向量 [RJRC25, Def. 3.2.3]。该函子是 Schneider--Teitelbaum 那个把 $G$ 的容许 Banach 空间表示送到其局部解析向量的函子的（导出）推广。于是 Rodrigues Jacinto 与 Rodríguez Camargo 把 $D(\mathrm{an}\,G)$ 定义为 $D_{L^\blacksquare}(\mathcal D(G))$ 的全子范畴，由满足「自然映射 $\pi^{\mathrm{Rla}}\to\pi$ 是同构」的复形 $\pi$ 组成 [RJRC25, Def. 3.3.1]。

\ysrc{256}
由于「导出局部解析」可以借非导出的局部解析向量在上同调群上检验 [RJRC25, Theorem A (1)]，范畴 $D(\mathrm{an}\,G)$ 的天然 $t$-结构之心便包含（非导出）局部解析表示所构成的范畴 $\mathrm{Rep}^{\mathrm{an}}_L G$。应用同一结果还可看出，$D(\mathrm{an}\,G)$ 包含 $D(\mathrm{Rep}^{\mathrm{sm}}G)$，即光滑 $G$-表示范畴的导出范畴（但并不作为全子范畴，因为 $\Ext$ 群并不相同）。

为了讨论从 $D(\mathrm{an}\,G)$ 到叠 $\mathcal X_d$（及其各种变体）上凝聚层的函子\tn{底本在 §6.2 各处时或把该叠印作 $\mathfrak X_d$（花体 X），而 §6.1 与本书其余部分印作 $\mathcal X_d$。译文统一采用 $\mathcal X_d$。}，我们须对 $D(\mathrm{an}\,G)$ 施加某个有限性条件，才能映到凝聚层（而不是固态拟凝聚层或 ind-凝聚层，这些我们不想在此讨论）。下文中我们因而写作
\[
D^b_{\mathrm{f.p.}}(\mathrm{an}\,G)\subset D(\mathrm{an}\,G)
\]
表示那些预期满足该条件的（导出）局部解析表示所构成的全子范畴。第一个猜测或许是：$D^b_{\mathrm{f.p.}}(\mathrm{an}\,G)$ 恰好由 $D(\mathrm{an}\,G)$ 中的紧对象组成；但有一些理由预期并非如此。正确的猜测似乎是：$D^b_{\mathrm{f.p.}}(\mathrm{an}\,G)$ 由那些关于局部解析表示上的六函子形式为 \emph{prim} 的对象所构成的全子范畴\footnote{不过读者也不妨近似地把 $D(\mathrm{an}\,G)$ 中的紧对象范畴当作它。}\tn{\emph{prim} 是原文保留的术语，此处指「相对于六函子形式不可再分」的对象；上下文推断，无外部依据。}。该范畴包含有限表现上同调、取值于 $L$-向量空间的光滑 $G$-表示的有界复形所构成的范畴 $D^b_{\mathrm{f.p.}}(\mathrm{Rep}^{\mathrm{sm}}G)$（同样并不作为全子范畴）。

\begin{rem}{6.2.2}
记号 $D^b_{\mathrm{f.p.}}(\mathrm{an}\,G)$ 相当不规范，但我们希望所用记号能让人联想到 Banach 情形中的记号。
\end{rem}

\subsubsection*{6.2.3.\quad 猜想的表述}

回忆一下，$\Q_p$-解析群 $G$ 的李代数 $\mathfrak g_L$ 的包络代数 $U(\mathfrak g_L)$ 的中心 $Z(\mathfrak g_L)$ 嵌入 $\mathcal D(G)$，[ST02, Prop. 3.7]。借助 Harish--Chandra 同构，该中心可以等同于 $\mathcal X_d$ 的 Hodge--Tate--Sen 权空间 $\mathrm{WT}_{d,L}$ 的整体截面 $\Gamma(\mathrm{WT}_{d,L},\mathcal O_{\mathrm{WT}_{d,L}})$。\footnote{回忆一下，本节中群 $G$ 就是 $\mathrm{GL}_d$，因此很容易把包络代数的中心与 Hodge--Tate--Sen 权空间上的函数空间作比较。对一般的约化群，这一比较更难，需要依赖与 [DPS25, 4.6] 中类似的构造。但无论如何，我们并未对任意约化群定义带 $G$-结构的 $(\varphi,\Gamma)$-模空间。}

进一步回忆 [BZCHN24, Theorem 5.13]、[Zhu25a, Conjecture 4.5.1]、[Hel23, Conjecture 3.6]（另见 §8.3）：对光滑表示，我们有一个函子
\[
\mathfrak A^{\mathrm{sm}}_G:D^b_{\mathrm{f.p.}}(\mathrm{Rep}^{\mathrm{sm}}G)\longrightarrow D^b_{\mathrm{coh}}(\mathrm{WD}_{d,F}).
\]
（对一般的约化群，该函子自然是猜想的；但对 $G=\mathrm{GL}_d(F)$，可以借型理论从 Iwahori--Hecke 代数的情形构造出该函子，见 [BZCHN24]。）

下面我们表述一个关于局部解析表示与 $\mathcal X_d$ 上凝聚层的猜想，它大致平行于光滑模 $p$ 表示与 $\mathcal X_d$ 上的层这一情形下的猜想 6.1.15。

\ysrc{257}
\begin{conj}{6.2.4}\label{conj:a4-6-2-4}
存在一个正合的 $\Z(\mathfrak g_L)$-线性函子
\[
\mathfrak A^{\mathrm{rig}}_G:D^b_{\mathrm{f.p.}}(\mathrm{an}\,G)\to D^b_{\mathrm{coh}}(\mathcal X_d),
\]
满足以下条件：

(1)（与光滑情形的相容性）设 $\xi$ 是对应于正则 Hodge--Tate 权 $\underline\lambda$ 的一个支配代数特征，并照定义 2.3.2 用 $W_{\underline\lambda}[1/p]$ 表示 $G$ 的相应的、最高权为 $\xi$ 的不可约代数表示。

设 $\mathrm{Rep}^{\mathrm{sm}}_{\Omega^\tau}G$ 是对应于惯性型 $\tau$ 的 Bernstein 块，其中 $\tau|_{I_{F'}}$ 对 $F$ 的某个有限扩张 $F'$ 平凡。经基变换到 $L'\supset F'$ 之后，有自然同构
\begin{equation}\tag{6.2.5}\label{eq:a4-6-2-5}
\mathfrak A^{\mathrm{rig}}_G(W_{\underline\lambda}[1/p]\otimes-)
\cong\mathrm{pr}^*_{\mathrm{WD}}\bigl(\mathfrak A^{\mathrm{sm}}_G(-)\bigr)
\bigl((-\rho')_{\sigma:F\into L}\bigr)
\end{equation}
其中
\[
\mathrm{pr}_{\mathrm{WD}}:\mathrm{Fil}_{\underline\lambda}\mathrm{Mod}_{d,\varphi,N,\tau}
\cong\mathrm{Fil}_{\underline\lambda}\mathrm{WD}_{d,F,\tau}
\longrightarrow\mathrm{WD}_{d,F,\tau}
\]
是典范投影。

(2)（与抛物诱导的相容性）设 $P\subseteq G$ 是抛物子群，其 Levi 商为 $M$，并用 $\hat P$、$\hat M$ 分别表示相应的对偶群。于是局部解析抛物诱导
\[
\mathrm{Ind}^G_{\bar P}(-)^{\mathrm{an}}:D(\mathrm{an}\,M)\to D(\mathrm{an}\,G),
\]
其中 $\bar P\subset G$ 是反抛物子群，满足相容性
\begin{equation}\tag{6.2.6}\label{eq:a4-6-2-6}
\mathfrak A^{\mathrm{rig}}_G\bigl(\mathrm{Ind}^G_{\bar P}(-)^{\mathrm{an}}\bigr)
\cong\tilde\beta_{\hat P,*}\circ\tilde\alpha^*_{\hat P}
\bigl(\mathfrak A^{\mathrm{rig}}_M(-)\bigl([F:\Q_p]\rho'_M\bigr)\bigr),
\end{equation}
只要该公式有意义；更精确的表述见 §6.2.14。
\end{conj}

\begin{rem}{6.2.7}
我们还预期有与对偶的相容性，类似于猜想 6.1.15 的 (3)。这会涉及在 $D(\mathrm{an}\,G)$ 上定义 Bernstein--Zelevinsky 型的对偶。应当可以通过把 $D(\mathrm{an}\,G)$ 解释为某个解析叠上拟凝聚层的范畴，并对解析叠使用六函子形式体系，来定义这样的对偶。不过我们在本讲义中不朝这个方向展开。
\end{rem}

\begin{rem}{6.2.8}
下面我们对 (6.2.5) 与 (6.2.6) 中分别出现的、以 $-\rho'$ 与以 $\rho'_M$ 所作的扭略作说明。

(i) 照惯例，用 $\rho$ 表示对偶环面 $\hat T$ 中正根之和的一半。为固定记号，取 $\hat T$ 为 $\mathrm{GL}_d$ 中的对角环面，并典范地等同 $\hat T=\mathbf G_m^d$。则
\[
\rho=\Bigl(\frac{d-1}{2},\frac{d-3}{2},\cdots,\frac{-d+3}{2},\frac{-d+1}{2}\Bigr)
\]
一般说来并不是 $\hat T$ 的特征，而只落在 $X^*(\hat T)\otimes\Q$ 中。我们改而考虑特征 $\rho'=(0,-1,\dots,-(d-1))\in X^*(\hat T)$。注意 $\rho'$ 是 $\rho$ 经 $X^*(\hat T)\otimes\Q$ 中某个元素平移所得，该元素经由行列式分解，且在 $\hat T$ 的余根上取值时与 $\rho$ 一致。

设 $M$ 是 $G$ 的 Levi 子群，则用 $\rho_M$ 表示在 $\mathrm{Lie}\,\bar U$ 上出现的那些根之和的一半，其中 $\bar U\subset\bar P$ 是幂幺根。于是同样可以对 $\rho_M$ 作平移，得到一个 $\hat M=\prod_{i=1}^r\mathrm{GL}_{r_i}$ 的特征 $\rho'_M$；\ysrc{258}它经由极大环面商 $S_{\hat M}=\prod_{i=1}^r\mathbf G_m$ 分解，并且 $\rho'_M$ 在 $S_{\hat M}$ 的第一个 $\mathbf G_m$ 因子上平凡。

(ii) 现在描述 (6.2.5) 中的扭。空间 $\mathrm{Fil}_{\underline\lambda}\mathrm{WD}_{d,F,\tau}$ 有典范投影到
\[
\mathrm{GL}_d\backslash(\mathrm{Res}_{F/\Q_p}\mathrm{GL}_d)_L/P_{\underline\lambda}
\cong\mathrm{GL}_d\backslash\Bigl(\prod_{\sigma:F\into L}\mathrm{GL}_d/\hat B\Bigr).
\]
这里 $P_{\underline\lambda}$ 表示 §5.2 中定义的那个抛物子群，由于 $\underline\lambda$ 正则，它与 Borel 的 Weil 限制一致。于是我们可以用旗簇之积上某个线丛的拉回来扭 $\mathrm{Fil}_{\underline\lambda}\mathrm{WD}_{d,F,\tau}$ 上的层，该线丛取为每个因子上与 $\hat T$ 的特征 $-\rho'$ 相伴的、$\mathrm{GL}_d/\hat B$ 上的线丛。我们要指出，以 $(-\rho')_\sigma$ 作扭与对偶是相容的（这与注 6.1.33 中的相容性精神一致）。同样，我们有 $D^b_{\mathrm{coh}}(\mathcal X_d)$ 上的自对偶 $\mathcal D_{\mathcal X_d}$，它由 Grothendieck--Serre 对偶与「$\mathrm{GL}_d$ 上主对合所诱导的映射 $\mathcal X_d\to\mathcal X_d$ 的拉回」复合而成。若 $\mathcal F$ 表示 $\mathrm{WD}_{d,F,\tau}$ 上的一个层复形，则
\[
\mathcal D_{\mathcal X_d}\bigl(\mathrm{pr}^{\mathrm{WD}}_{\mathrm{WD}}(\mathcal F)(-\rho')_\sigma\bigr)
=\mathrm{pr}^{\mathrm{WD}}_{\mathrm{WD}}\bigl(\mathcal D_{\mathrm{WD}_{d,F,\tau}}(\mathcal F)\bigr)(-\rho')_\sigma,
\]
其中 $\mathcal D_{\mathrm{WD}_{d,F,\tau}}$ 的定义方式与 $\mathcal D_{\mathcal X_d}$ 相同\tn{底本此处 $\mathrm{pr}$ 的上标与下标都印作 $\mathrm{WD}$（对比 (6.2.5) 中的 $\mathrm{pr}^*_{\mathrm{WD}}$），疑上标当作 $*$。译文照底本原样保留。}。

(iii) 最后描述 (6.2.6) 中的扭。行列式诱导一个典范态射 $\mathcal X_d\to\mathcal X_1\to*/\mathbf G_m$（关于 $\mathcal X_1$ 的描述以及到 $*/\mathbf G_m$ 的典范投影，见 (7.1.3)）。特别地，我们得到一个投影
\[
\mathcal X_{\hat M}\to*/S_{\hat M},
\]
于是用 $S_{\hat M}$ 的特征来扭 $\mathcal X_{\hat M}$ 上的层是有意义的（我们把该特征看作 $*/S_{\hat M}$ 上的线丛，并把它拉回到 $\mathcal X_{\hat M}$）。

还应当有进一步的相容性，即与下文 §6.2.11 所讨论的 Banach 情形的相容性，以及我们将在 §9.6 讨论的局部--整体相容性。下面我们对猜想中的这些相容性略作评论。
\end{rem}

\begin{rem}{6.2.9}
(a) 性质 (1) 把函子 $\mathfrak A^{\mathrm{rig}}_G$ 在局部代数表示上确定下来。这意味着，我们预期：把光滑表示上的函子 $\mathfrak A^{\mathrm{sm}}_G$ 扩充到局部代数表示所得者（它取值于支撑在 de Rham 轨迹上的凝聚层）可以在某种意义上插值成一个取值于 $\mathcal X_d$ 上凝聚层（对支撑不加限制）的函子。这在某种意义上平行于下述观点：$p$-进自守形式（上的 Hecke 特征值）插值经典自守形式（上的 Hecke 特征值）。

(b) 性质 (1) 中的自然同构依赖于引理 5.2.3 中同构
\[
\mathrm{Fil}_{\underline\lambda}\mathrm{Mod}_{d,\varphi,N,\tau}
\cong\mathrm{Fil}_{\underline\lambda}\mathrm{WD}_{d,F,\tau}
\]
的选取，即依赖于嵌入 $F'\into L'$ 的选取。事实上我们预期函子 $\mathfrak A^{\mathrm{sm}}_G$ 有一个变体，它是函子
\[
D^b_{\mathrm{f.p.}}(\mathrm{Rep}^{\mathrm{sm}}G)
\longrightarrow D^b_{\mathrm{coh}}\Bigl(\coprod_{F'/F}\mathrm{Mod}_{d,\varphi,N,F'/F}\Bigr),
\]
并在作适当的扩域之后与 $\mathfrak A^{\mathrm{sm}}_G$ 一致，见注 8.3.3。这一版本的 $\mathfrak A^{\mathrm{sm}}_G$ 将使性质 (1) 的表述更顺畅（即无需辅助地扩域到一个包含 $F'$ 的域）。
\end{rem}

\ysrc{259}
(c) 我们无法指望函子 $\mathfrak A^{\mathrm{rig}}_G$ 在不对来源 $D^b_{\mathrm{f.p.}}(\mathrm{an}\,G)$ 施加任何额外条件的情况下就是全忠实的。在 $\mathrm{GL}_1$ 情形，这一点在下文注 7.1.11 中讨论。但这一现象在局部代数表示上也能看到。函子 $\mathfrak A^{\mathrm{sm}}_G$ 是（猜想地）全忠实的。给定两个光滑表示 $\pi_1$ 与 $\pi_2$（集中在 0 次），同态模 $\Hom_G(\pi_1,\pi_2)$ 无论在光滑表示范畴中计算还是在局部解析表示范畴中计算都是同一个（高阶 $\Ext$-群当然不是这样）。由全忠实性，该模等于 $\Hom\bigl(\mathfrak A^{\mathrm{sm}}_G(\pi_1),\mathfrak A^{\mathrm{sm}}_G(\pi_2)\bigr)$；但一般说来，这些层的解析化之间的同态空间严格更大。看来只要再加上一个\emph{有限斜率}条件，或要求所论表示容许，就不会出现这一问题；可与 $\mathrm{GL}_1$ 情形中的讨论相比较。遗憾的是，同时施加「在 $\mathcal D(G)$ 上有限表现」与「容许」这两个假设是相当强的限制，会把许多有意思的表示排除在外。

(d) 我们预期，若允许比凝聚层更一般的（复形的）层，则函子 $\mathfrak A^{\mathrm{rig}}_G$ 可以定义在像 $D(\mathrm{an}\,G)$ 这样更大的范畴上，而不必限制在 $D^b_{\mathrm{f.p.}}(\mathrm{an}\,G)$ 上。与 $p$-幂挠表示的情形类似，我们预期此时表示 $\pi$ 的容许性并不保证 $\mathfrak A^{\mathrm{rig}}_G(\pi)$ 是凝聚层，除非对所论局部解析表示作为 $\mathcal D(G)$-模再施加某种额外的有限性假设。在 $p$-挠情形，这一点在 §7.7.59 中讨论；并且已知 [Paš10, Thm. 1.1]：当 $G=\mathrm{GL}_2(F)$ 时，局部代数 $G$-表示的酉 Banach 完备化，其约化可以包含任意给定的不可约容许模 $p$ 表示。因此我们预期那一讨论中的模 $p$ 现象在特征 0 中也有类似物。

\begin{rem}{6.2.10}
我们预期，与注 6.1.24 类似，猜想 6.2.4 中的函子 $\mathfrak A^{\mathrm{rig}}_G$ 应由
\[
\pi\longmapsto L_\infty\widehat{\otimes}^L_{\mathcal D(G)}\pi
\]
给出，其中 $-\widehat{\otimes}^L_{\mathcal D(G)}-$ 是某种导出完备张量积，而 $\mathcal L_\infty$ 是叠 $\mathcal X_d$ 上的一族 $\mathcal D(G)$-模。这一\emph{导出完备}张量积当然只在固态模与固态局部解析表示的世界中才有意义。我们预期，借助 (5.1.31) 中的映射 $\pi_d$，$\mathcal L_\infty$ 到 $\mathcal X_d$ 的一般纤维上的拉回应由
\[
\pi_d^*\mathcal L_\infty\cong L_\infty\otimes_{\mathcal O_{\mathcal X_d}[[G]]}\mathcal D(G,\mathcal O_{\mathcal X_d})
\]
给出，其中 $\mathcal D(G,\mathcal O_{\mathcal X_d})$ 是分布代数的一个相对版本，其定义方式与 [BHS17b, 3.1] 或 [GRJ18, 1.3] 中类似。

当 $G=\mathrm{GL}_2(\Q_p)$ 时，我们预期：对由 $(\varphi,\Gamma)$-模 $D$ 给出的点 $x\in\mathcal X_2(L)$，$\mathcal D(G)$-模 $\mathcal L_\infty\otimes k(x)$ 就是 Colmez 把 $p$-进 Langlands 对应推广到 $\mathcal R_F$ 上（不必 étale 的）$(\varphi,\Gamma)$-模情形之后，与 $D$ 相伴的局部解析表示的对偶 [Col16]。更确切地说：族 $\mathcal L_\infty$ 应当插值 Colmez 关于 $\mathcal R_L$ 上 $(\varphi,\Gamma)$-模的 $p$-进 Langlands 对应。利用注 6.1.24 中的对象 $L_\infty$，我们可以在 $\mathcal X_2$ 中「$(\varphi,\Gamma)$-模在扭的意义下 étale」的开子叠上（在 $\mathrm{GL}_2(\Q_p)$ 的情形）构造 $\mathcal L_\infty$。该开子集的补集只由三角化 $(\varphi,\Gamma)$-模组成。另一方面，\ysrc{260}Gaisin 与 Rodrigues Jacinto [GRJ18] 在参数化三角化 $(\varphi,\Gamma)$-模（连同一个三角化结构）的叠 $\mathcal X_B$ 的（正则部分）上，构造了 $\mathcal L_\infty$ 的一个类似物。至少目前还不清楚如何由这两种情形粘出 $\mathcal L_\infty$。
\end{rem}

\subsubsection*{6.2.11.\quad 与 Banach 情形的相容性}

我们预期 Banach 情形下的猜想（猜想 6.1.15）与解析情形下的猜想（猜想 6.2.4）之间有很强的联系。在表示论一侧，这一相容性应当涉及从 Banach 空间表示 $V=\Lambda[1/p]$ 中的格 $\Lambda$（把它看作 $D^b_{\mathrm{f.p.}}(\mathrm{sm}\,G)$ 中的 pro-对象）过渡到 $V$ 中的局部解析向量（参看 [ST03, §7]）。在 $(\varphi,\Gamma)$-模叠上的（凝聚）层这一侧，我们考虑 $D^b_{\mathrm{coh}}(\mathcal X_d)$ 中某个 pro-对象的一般纤维 $\mathcal F_\eta$。它（至少在启发式上如此；关于 $\mathrm{GL}_1$ 情形的讨论见引理 7.1.9 之后的讨论）给出 $\mathcal X^{\mathrm{rig}}_{d,\eta}$ 上的一个层复形，而我们把这个层复形与「沿着 (5.1.31) 对 $\mathcal X_d$ 上某个层复形的拉回」作比较。更确切地说，对 pro-对象 $\Lambda\in D^b_{\mathrm{f.p.}}(\mathrm{sm}\,G)$ 且 $V=\Lambda[1/p]$，该比较应当形如
\begin{equation}\tag{6.2.12}\label{eq:a4-6-2-12}
\mathfrak A(\Lambda)^{\mathrm{rig}}_\eta\cong\pi_d^*\mathfrak A^{\mathrm{rig}}_G(V^{\mathrm{an}}).
\end{equation}

由于我们所施加的那些有限性条件，这一公式会遇到一些困难。我们预期，只要在两侧都使用固态模（即固态局部解析表示与固态拟凝聚层），这些困难便可处理；但我们不在此朝这个方向展开。下面我们改为讨论人们必须处理的一些困难。

\begin{itemize}
\item 我们一直局限于刚性解析空间（以及刚性解析 Artin 叠）上的凝聚层，尽管这些空间上应当有一个更灵活的环境范畴（用固态模来定义），可与注 5.1.11 相比较。虽然可以对 $\mathcal X_d$ 上的 pro-凝聚层 $\mathcal F$ 施加条件，以保证 $\mathcal F$ 有良定义的一般纤维 $\mathcal F^{\mathrm{rig}}_\eta$，且后者是 $\mathcal X^{\mathrm{rig}}_{d,\eta}$ 中的凝聚层，但不太清楚应当对 $D(\mathrm{sm}\,G)$ 中的 pro-对象 $\Lambda$ 施加哪些有限性条件，才能保证 $\mathfrak A(\Lambda)^{\mathrm{rig}}_\eta$ 是 $D_{\mathrm{coh}}(\mathcal X^{\mathrm{rig}}_{d,\eta})$ 的对象。给定单个表示 $\Lambda$（比如处于上同调次数 0），该有限性假设至少应涉及在 $\mathcal O[[G]]$ 上的有限性。

\item 在我们的导出范畴语境中，过渡到局部解析向量 $V\mapsto V^{\mathrm{an}}$ 应当是导出的。结果表明，在容许 Banach 空间表示的范畴上这并不必要。更确切地说，若 $V$ 是容许 Banach 空间表示（即其对偶在 $L[[K]]$ 上有限生成），则 $V^{\mathrm{an}}$ 是容许局部解析表示，且其对偶可以如下描述：由 [ST03, §4]，存在典范平坦映射 $\mathcal O[[K]]\to\mathcal D(K)$，它从完备群环 $\mathcal O[[K]]$ 到分布代数 $\mathcal D(K)$。给定有限生成 $\mathcal O[[K]]$-模 $\pi$，可以考虑标量扩张 $\pi\otimes_{\mathcal O[[K]]}\mathcal D(K)$。若 $\pi$ 带有 $G$-作用（与 $\mathcal O[[K]]$-模结构相容），则 $\pi\otimes_{\mathcal O[[K]]}\mathcal D(K)$ 自然地成为余容许 $\mathcal D(G)$-模。由于 $\mathcal O[[K]]\to\mathcal D(K)$ 平坦，函子 $(-)^{\mathrm{an}}$ 在容许 Banach 空间表示的范畴上正合；并且事实上可以推出 $V^{\mathrm{an}}$ 甚至是导出局部解析的，见 [RJRC22, Prop. 4.48]、[Pan22b, Theorem 2.2.3]。
\ysrc{261}
\item 正如注 6.2.9 所论，在处理函子 $\mathfrak A$ 与 $\mathfrak A^{\mathrm{rig}}_G$ 时，容许性条件并不是一个好的有限性条件：一般说来，我们不应预期容许表示会被送到凝聚层。另一方面，同时施加两个有限性条件——容许性，以及在 $\mathcal O[[G]]$（相应地 $\mathcal D(G)$）上的有限性——限制相当强，会把许多有意思的表示排除在外。因此看来，只有借助固态模理论扩大我们的环境范畴，才能给出相容性 (6.2.12) 的一般表述。
\end{itemize}

Banach 情形与解析情形的相容性 (6.2.12) 应当有以下一个实例。给定正则 Hodge--Tate 权 $\underline\lambda$ 与惯性型 $\tau$，猜想 6.1.15 (4) 预测
\[
\mathcal F:=\mathfrak A\bigl(\bigl(\widehat{\Pi(\underline\lambda,\tau)}\bigr)^{\circ}\bigr)
\]
是集中在 0 次、支撑在 $p$-进叠 $\mathcal X_d^{\mathrm{crys},\underline\lambda,\tau}$ 上的 pro-凝聚层。我们预期，对任意态射 $\Spf A\to\mathcal X_d^{\mathrm{crys},\underline\lambda,\tau}$（其中 $A$ 是 $p$-进完备、在 $\mathcal O$ 上拓扑有限型的 $\mathcal O$-代数），它的拉回由有限生成 $A$-模给出。此时 $\mathcal F$ 有良定义的一般纤维 $\mathcal F^{\mathrm{rig}}_\eta$，它是 $(\mathcal X_d^{\mathrm{crys},\underline\lambda,\tau})^{\mathrm{rig}}_\eta$ 上的凝聚层。于是我们预期有同构
\begin{equation}\tag{6.2.13}\label{eq:a4-6-2-13}
\mathcal F^{\mathrm{rig}}_\eta\cong\pi_d^*\bigl(\mathfrak A^{\mathrm{rig}}_G(\Pi(\underline\lambda,\tau))\bigr).
\end{equation}

正如已经指出的，这大概是猜想性相容性 (6.2.12) 的一个实例；但要想更充分地理解它究竟如何实现，似乎需要描述 $\Pi(\underline\lambda,\tau)^\circ$ 中局部解析向量的子空间。给出这样的描述本身就是一个有意思的问题，遗憾的是我们目前还不知道答案。

最后我们要指出，Colmez [Col19, Rem. 0.2] 描述了：对秩 2 的 étale $(\varphi,\Gamma)$-模 $D$，在 $\mathrm{GL}_2(\Q_p)$ 的 $p$-进局部 Langlands 对应下与之相伴的 Banach 空间表示 $\Pi$ 中，局部解析向量 $\Pi(D)^{\mathrm{an}}$ 的 Jordan--Hölder 因子。预测 [Col19, Corollaire 0.4] 中的那些表示在函子 $\mathfrak A^{\mathrm{rig}}_G$ 下被送到何处，将会很有意思。在三角化情形，这基本上已由所预期的「与光滑情形的相容性」和「与抛物诱导的相容性」涵盖。但对剩下的表示，我们不知道如何（猜想地）描述其像。

\subsubsection*{6.2.14.\quad 与抛物诱导的相容性}

固定一个具有分裂极大环面 $T$ 的 Borel 子群 $B\subseteq G$。再设 $P\supseteq B$ 是 Levi 为 $M$ 的抛物子群。设 $(\pi,V)$ 是 $M$ 在紧型 $L$-向量空间 $V$ 上的局部解析表示。$\pi$ 的局部解析抛物诱导定义为
\[
(\mathrm{Ind}^G_P\pi)^{\mathrm{an}}
=\{f:G\to V\ \text{局部解析}\mid f(gp)=\pi(p)^{-1}f(g),\ 
\text{对一切 }p\in P,\ g\in G\}.
\]
特别地，我们得到函子
\begin{align*}
\mathrm{Ind}^G_P(-)^{\mathrm{an}}&:\mathrm{Rep}^{\mathrm{an}}_L M\longrightarrow\mathrm{Rep}^{\mathrm{an}}_L G,\\
\mathrm{Ind}^G_P(-)^{\mathrm{an}}&:D(\mathrm{an}\,M)\longrightarrow D(\mathrm{an}\,G),
\end{align*}

\ysrc{262}
其中我们滥用记号，对（非形式地定义的）导出范畴上的函子使用同一记号。为使讨论不致过于繁复，下文只关注 $P=B$ 的情形。记
\[
s_G:D(\mathrm{an}\,T)\longrightarrow D(\mathrm{an}\,T)
\]
为用
\begin{equation}\tag{6.2.15}\label{eq:a4-6-2-15}
\delta_B\cdot\bigl(1,\dots,(\varepsilon\circ\mathrm{rec})^{i-1},\dots,
(\varepsilon\circ\mathrm{rec})^{d-1}\bigr)
\end{equation}
扭 $T$-表示的自函子\footnote{我们把这个平移与扭按 $G$ 而不按 $T$ 编号，因为该平移与扭显然依赖于 $G$，而不只依赖于 $T$。}（其中 $\delta_B$ 是 $B$ 的光滑模特征，$\varepsilon$ 是分圆特征，$\mathrm{rec}$ 是局部类域论的同构）。

\medskip\noindent\textbf{期望 6.2.16.}\label{exp:a4-6-2-16} 存在同构
\[
\mathfrak A^{\mathrm{rig}}_G\bigl(\mathrm{Ind}^G_B(\pi)^{\mathrm{an}}\bigr)
\cong\tilde\beta_{\hat B,*}\tilde\alpha^*_{\hat B}
\bigl(\mathfrak A^{\mathrm{rig}}_T\bigl(s_G(\pi)\bigl[[F:\Q_p]\rho'\bigr]\bigr)\bigr),
\]
对 $\pi$ 函子化，只要 $\pi$ 是使得 $\mathrm{Ind}^G_B(\pi)^{\mathrm{an}}$ 落在 $D^b_{\mathrm{f.p.}}(\mathrm{an}\,G)$ 中、且 $\tilde\alpha^*_{\hat B}\bigl(\mathfrak A^{\mathrm{rig}}_T(\pi)\bigr)$ 的支撑在 $\mathcal X_d$ 上是真（proper）的表示。

\begin{rem}{6.2.17}
以特征 (6.2.15) 作扭这一现象，是由所选定的正规化造成的。事实上，我们所选的正规化依赖于函子 $\mathfrak A^{\mathrm{sm}}_G$ 的正规化，或者（二者大致相同）依赖于命题 5.2.12 中同构的正规化，可与注 5.2.13 相比较。事实上我们这里所用的正规化与 [Hel23] 中所用的不同，那里与抛物诱导的相容性是在不涉及扭的情况下陈述的（但用的是正规化抛物诱导）。[Hel23] 中的正规化对应于局部 Langlands 对应的酉正规化，而不是我们这里所用的非酉正规化。为此付出的代价，就是期望 6.2.16 中出现了这个扭。
\end{rem}

我们要强调，不应指望局部解析抛物诱导保持有限性条件，即对 $\pi\in D^b_{\mathrm{f.p.}}(\mathrm{an}\,T)$，我们并不预期 $\mathrm{Ind}^G_B(\pi)^{\mathrm{an}}$ 落在 $D^b_{\mathrm{f.p.}}(\mathrm{an}\,G)$ 中。不过若 $\pi$ 具有有界斜率，这应当是对的。我们预期这一现象与「$\tilde\alpha^*_{\hat B}\bigl(\mathfrak A^{\mathrm{rig}}_T(\pi)\bigr)$ 的支撑在 $\mathcal X_d$ 上是否为真」这一问题直接相关。此外，正如上文已经指出的，我们预期存在 $\mathfrak A^{\mathrm{rig}}_G$ 的一个变体，它定义在 $D(\mathrm{an}\,G)$ 上，取值于 $\mathcal X_d$ 上一个更大的层范畴（最可能是固态拟凝聚层）。在这一语境中，人们应当对这些层有一个完整的六函子形式体系；此时与抛物诱导的相容性 6.2.16 更应当用 $\tilde\beta_{\hat B,!}$ 而非 $\tilde\beta_{\hat B,*}$ 来表述（当然，若假定相应层的支撑为真（proper），我们就看不出二者的差别）。

我们要提到，Orlik 与 Strauch [OS15] 描述了某些局部解析主级数表示的（拓扑）不可约子商。更确切地说，他们描述了 $\bigl(\mathrm{Ind}^G_B\delta\bigr)^{\mathrm{an}}$ 的不可约子商，其中 $\delta:T\to L^\times$ 是局部代数特征。他们研究的主要工具是构造某些双函子 $\mathcal F^G_P(-,-)$：它把一对 $(M,\pi)$ 送到一个

% 本片止于印刷页 262 页末。原文最后一句在此处被页末截断（"…to a"），续片 ch-a4-09 从 PDF 276 页（印刷页 263）起续译该句。
